\section{Research programme}\label{research-programme}

\subsection{1. The full structural chain}\label{the-full-structural-chain}

The model can be summarized as a composition of operators:

\[
\mu_t
\xrightarrow{\mathcal G}
\mu_{t+1}
\xrightarrow{\mathcal H,\mathcal B}
\widetilde\Lambda_t
\xrightarrow{\mathcal M}
\Gamma_t^*
\xrightarrow{\mathcal P}
\text{economic outcomes}
\xrightarrow{\mathcal O}
Y_{t+1}
\xrightarrow{\mathcal F}
\pi_{t+1}.
\]

Interpretation:

\begin{itemize}
\tightlist
\item
  \texttt{G}: stochastic growth, mortality, recruitment, and quality transition;
\item
  \texttt{H}: harvest selection;
\item
  \texttt{B}: bucking/recovery from tree to logs;
\item
  \texttt{M}: processor/operator matching;
\item
  \texttt{P}: price, surplus, and bargaining layer;
\item
  \texttt{O}: observation process;
\item
  \texttt{F}: Bayesian filtering/update.
\end{itemize}

The controller chooses actions that can influence several of these operators at
once.

This operator view is the most reusable part of the theory. Each operator can be
implemented first as a coarse approximation and later replaced by a richer,
identified model while preserving the interfaces around it.

\subsection{2. Why separation creates falsifiability}\label{why-separation-creates-falsifiability}

Suppose a predicted forest-gate price is wrong. A scalar model may provide no
clear diagnosis. The structural model asks where the error entered:

\begin{itemize}
\tightlist
\item
  Was standing size/quality wrong?
\item
  Was tree-to-log recovery wrong?
\item
  Was processor compatibility wrong?
\item
  Was haulage or operator productivity wrong?
\item
  Was the outside option wrong?
\item
  Was bargaining behavior wrong?
\end{itemize}

An accurate EO-derived biomass estimate can coexist with a poor commercial-supply
prediction when a downstream operator is misspecified.

This decomposition turns model error into research questions.

\subsection{3. Mathematical open problems}\label{mathematical-open-problems}

\subsubsection{3.1 Joint finite-mass dynamics with common noise}\label{joint-finite-mass-dynamics-with-common-noise}

The paper combines a normalized conditional McKean-Vlasov law with a
mass-changing finite measure. A deeper theory could formulate the full
stochastic transport-diffusion-reaction process directly on an HK-like
finite-measure space with common noise, controls, and killing/recruitment.

Questions include existence, stability, propagation of chaos for weighted
particles, and a useful dynamic-programming principle on the resulting space.

\subsubsection{3.2 Stratified geometry}\label{stratified-geometry}

Species, genetic classes, and regimes are discrete; size, location, and quality
are continuous. The global state is a stratified space.

An open problem is to characterize optimization and viscosity solutions on this
hybrid geometry without pretending that discrete class changes are continuous
transport.

\subsubsection{3.3 Measure-valued impulse control}\label{measure-valued-impulse-control}

Harvest and regime changes can be impulsive. A rigorous quasi-variational
inequality on a belief space whose physical component contains finite measures
would connect the paper's formal recursion to a stronger analytical theory.

\subsubsection{3.4 Endogenous bucking and matching}\label{endogenous-bucking-and-matching}

When the bucking value function is taken from downstream matching duals, physical
recovery and market allocation become coupled. Conditions for fixed points,
stability, or decomposable solution methods are natural theoretical questions.

\subsection{4. Statistical open problems}\label{statistical-open-problems}

\subsubsection{4.1 Growth diffusion versus heterogeneity}\label{growth-diffusion-versus-heterogeneity}

Repeated-tree data are needed to distinguish process noise, measurement error,
site effects, genetic effects, and latent micro-site variation. Cross-sectional
spread alone supplies no estimate of process diffusion.

\subsubsection{4.2 Competition identification}\label{competition-identification}

The structural model allows a general law-dependent interaction. Which summary
or kernel is identifiable from East African trial and plantation data is an
empirical question.

Candidates include:

\begin{itemize}
\tightlist
\item
  stem density;
\item
  basal area;
\item
  stand-density index;
\item
  size-asymmetric crowding;
\item
  spatial neighbour kernels.
\end{itemize}

Out-of-site predictive performance and biological interpretability determine the
appropriate level of complexity.

\subsubsection{4.3 Recovery and grade transition}\label{recovery-and-grade-transition}

The tree-to-log kernel is a major missing link. It requires pre-harvest tree
state joined to log dimensions, defect, moisture/density, grade, and processor
acceptance.

\subsubsection{4.4 Observation discrepancy}\label{observation-discrepancy}

EO, field plots, management records, and processor intake observe different
functions of the latent state. Cross-sensor reconciliation and structural
model-discrepancy modelling are open statistical tasks.

\subsection{5. Economic and market open problems}\label{economic-and-market-open-problems}

\subsubsection{5.1 Processor demand response}\label{processor-demand-response}

Posted prices and marginal willingness to pay can differ. Quantity response,
minimum lots, stockout risk, and production schedules can make demand
state-dependent and non-convex.

\subsubsection{5.2 Outside options and bargaining}\label{outside-options-and-bargaining}

The stylized Nash split is a placeholder until repeated transactions, bids, and
alternative-buyer observations identify a stable bargaining structure.

\subsubsection{5.3 Network effects and adoption}\label{network-effects-and-adoption}

A two-sided forestry intelligence platform may become more valuable as more
asset owners, processors, and operators participate. That can destroy the simple
one-segment log-concavity result and create multiple equilibria or threshold
behavior.

\subsection{6. Decision-science open problems}\label{decision-science-open-problems}

\subsubsection{6.1 Which measurements are worth buying?}\label{which-measurements-are-worth-buying}

The posterior framework allows a direct question: where does one additional
field plot, processor intake measurement, or contractor observation have the
highest expected decision value?

The matching dual suggests a powerful heuristic: prioritize uncertainty in
regions of state space with high marginal supply value and high classification
uncertainty.

\subsubsection{6.2 When should coverage scale?}\label{when-should-coverage-scale}

Regime switching depends on posterior evidence and the irreversibility of scale
costs. The practical research task is to define observable sufficient statistics
for a tractable control rule.

\subsubsection{6.3 How much model complexity pays for itself?}\label{how-much-model-complexity-pays-for-itself}

Model selection maximizes expected decision improvement net of sensing, compute,
and operating cost. This criterion can favor a simpler model with reliable input
data.

\subsection{7. A proposed research sequence}\label{a-proposed-research-sequence}

A sensible sequence implied by the paper is:

\begin{enumerate}
\def\labelenumi{\arabic{enumi}.}
\tightlist
\item
  establish stable canonical state and provenance contracts;
\item
  assemble repeated-tree and harvest-linked recovery data;
\item
  identify growth, mortality, and recovery models with explicit uncertainty;
\item
  build processor acceptance and delivered-cost observations;
\item
  benchmark the discrete matching LP and its duals;
\item
  connect high-value uncertainty to targeted field verification;
\item
  fit reduced belief coordinates for sequential decisions;
\item
  test regime-switching rules in retrospective or simulated experiments;
\item
  evaluate Bayes-adaptive operational performance against a fixed benchmark.
\end{enumerate}

\subsection{8. A theorem/proof agenda}\label{a-theoremproof-agenda}

The current working paper contains existence, convexity, continuity, and
structural propositions. A more mature mathematical paper could deepen the
proof programme around:

\begin{itemize}
\tightlist
\item
  convergence of weighted finite-particle approximations under killing and
  recruitment;
\item
  stability of the tree-to-log measure map with respect to perturbations in
  tree state and bucking kernels;
\item
  differentiability or subdifferential structure of matching value with respect
  to supply measures;
\item
  error bounds for entropy-regularized matching duals;
\item
  posterior consistency under designed observation policies;
\item
  conditions under which low-dimensional belief summaries are approximately
  sufficient for the outer control problem.
\end{itemize}

These should be treated as open until complete assumptions and proofs are
written.

\subsection{9. Product/research boundary}\label{productresearch-boundary}

The product should be allowed to use simpler models than the theory whenever the
simplification is explicit and empirically honest.

A production output might use a deterministic yield table while the research
branch evaluates stochastic size distributions. It might use observed processor
grade tables before a full quantity-response model exists. Scheduled verification
can approximate the EVSI policy.

The key standard is traceability:

\begin{Shaded}
\begin{Highlighting}[]
\NormalTok{production approximation}
\NormalTok{    {-}\textgreater{} structural operator it approximates}
\NormalTok{    {-}\textgreater{} evidence supporting the approximation}
\NormalTok{    {-}\textgreater{} known failure modes}
\NormalTok{    {-}\textgreater{} path to replacement}
\end{Highlighting}
\end{Shaded}

This traceability connects each production approximation to the research model.
